\nonstopmode{}
\documentclass[a4paper]{book}
\usepackage[times,inconsolata,hyper]{Rd}
\usepackage{makeidx}
\makeatletter\@ifl@t@r\fmtversion{2018/04/01}{}{\usepackage[utf8]{inputenc}}\makeatother
% \usepackage{graphicx} % @USE GRAPHICX@
\makeindex{}
\begin{document}
\chapter*{}
\begin{center}
{\textbf{\huge Package `PheValuator'}}
\par\bigskip{\large \today}
\end{center}
\ifthenelse{\boolean{Rd@use@hyper}}{\hypersetup{pdftitle = {PheValuator: Rapid Phenotype Evaluation}}}{}
\ifthenelse{\boolean{Rd@use@hyper}}{\hypersetup{pdfauthor = {William J. O'Brien; Joel Swerdel; Martijn Schuemie}}}{}
\begin{description}
\raggedright{}
\item[Type]\AsIs{Package}
\item[Title]\AsIs{Rapid Phenotype Evaluation}
\item[Version]\AsIs{2.2.19}
\item[Date]\AsIs{2026-10-7}
\item[Description]\AsIs{Evaluates phenotype algorithms by developing a diagnostic
prediction model using a set of noisy positives and noisy negatives.
It can then apply the model to a large cohort
to create a set of subjects each with a predicted probability for the
health outcome of interest.  This cohort can then be used to evaluate
phenotype algorithms through the calculation of the necessary performance
characteristics - sensitivity, specificity, and positive and negative
predictive value. Based on work by Swerdel (2022) <}\Rhref{https://doi.org/10.1016/j.jbi.2022.104177}{doi:10.1016/j.jbi.2022.104177}\AsIs{>
and Swerdel (2019) <}\Rhref{https://doi.org/10.1016/j.jbi.2019.103258}{doi:10.1016/j.jbi.2019.103258}\AsIs{>. }
\item[License]\AsIs{Apache License 2.0}
\item[VignetteBuilder]\AsIs{knitr}
\item[Depends]\AsIs{R (>= 4.1.1)}
\item[Imports]\AsIs{DatabaseConnector,
PatientLevelPrediction (>= 6.0.4),
FeatureExtraction (>= 2.1.5),
SqlRender (>= 1.5.2),
dplyr (>= 0.7.8),
methods,
ParallelLogger,
tibble}
\item[Suggests]\AsIs{knitr,
rmarkdown,
shiny,
shinydashboard,
DT,
htmltools,
qpdf,
testthat,
Eunomia,
withr,
CohortGenerator}
\item[URL]\AsIs{}\url{https://ohdsi.github.io/PheValuator/}\AsIs{, }\url{https://github.com/OHDSI/PheValuator}\AsIs{}
\item[BugReports]\AsIs{}\url{https://github.com/OHDSI/PheValuator/issues}\AsIs{}
\item[Encoding]\AsIs{UTF-8}
\item[Config/roxygen2/version]\AsIs{8.1.1}
\end{description}
\Rdcontents{Contents}
\HeaderA{computePerformanceMetricsFromCounts}{Compute performance metrics based on a 2-by-2 counts table}{computePerformanceMetricsFromCounts}
%
\begin{Description}
Compute performance metrics based on a 2-by-2 counts table
\end{Description}
%
\begin{Usage}
\begin{verbatim}
computePerformanceMetricsFromCounts(counts)
\end{verbatim}
\end{Usage}
%
\begin{Arguments}
\begin{ldescription}
\item[\code{counts}] Counts as created by the \code{\LinkA{testPhenotypeAlgorithm}{testPhenotypeAlgorithm}} function.
\end{ldescription}
\end{Arguments}
%
\begin{Value}
A tibble with the statistics metrics added to the counts table.
\end{Value}
%
\begin{Examples}
\begin{ExampleCode}
counts <- tibble::tibble(
  truePositives = 120,
  trueNegatives = 800,
  falsePositives = 30,
  falseNegatives = 50,
  cutPoint = "EV"
)
computePerformanceMetricsFromCounts(counts)

\end{ExampleCode}
\end{Examples}
\HeaderA{createCreateEvaluationCohortArgs}{Create a parameter object for the function createEvaluationCohort}{createCreateEvaluationCohortArgs}
%
\begin{Description}
Create a parameter object for the function createEvaluationCohort
\end{Description}
%
\begin{Usage}
\begin{verbatim}
createCreateEvaluationCohortArgs(
  phenotype,
  analysisName,
  runDateTime,
  databaseId,
  xSpecCohortId,
  daysFromxSpec = 0,
  xSensCohortId,
  prevalenceCohortId,
  caseCohortId,
  caseFirstOccurrenceOnly,
  xSpecCohortSize = 5000,
  covariateSettings = createDefaultCovariateSettings(excludedCovariateConceptIds = c(),
    addDescendantsToExclude = TRUE),
  modelPopulationCohortId = 0,
  modelPopulationCohortIdStartDay = 0,
  modelPopulationCohortIdEndDay = 0,
  inclusionEvaluationCohortId = 0,
  inclusionEvaluationDaysFromStart = 0,
  inclusionEvaluationDaysFromEnd = 0,
  duringInclusionEvaluationOnly = FALSE,
  exclusionEvaluationCohortId = 0,
  exclusionEvaluationDaysFromStart = 0,
  exclusionEvaluationDaysFromEnd = 0,
  priorModelToUse = NULL,
  minimumOffsetFromStart = 365,
  minimumOffsetFromEnd = 365,
  modelBaseSampleSize = 25000,
  baseSampleSize = 2e+06,
  lowerAgeLimit = 0,
  upperAgeLimit = 120,
  visitLength = 0,
  visitType = c(9201, 9202, 9203, 262, 581477),
  gender = c(8507, 8532),
  race = 0,
  ethnicity = 0,
  startDate = "19001010",
  endDate = "21000101",
  falsePositiveNegativeSubjects = 10,
  exportFolder,
  modelId = "main",
  excludeModelFromEvaluation = FALSE,
  randomVisitTable = "",
  removeSubjectsWithFutureDates = TRUE,
  saveEvaluationCohortPlpData = FALSE
)
\end{verbatim}
\end{Usage}
%
\begin{Arguments}
\begin{ldescription}
\item[\code{phenotype}] Name of the phenotype for analysis

\item[\code{analysisName}] Name of the analysis

\item[\code{runDateTime}] Starting date and time of the PheValuator run

\item[\code{databaseId}] Name of the database in the analysis

\item[\code{xSpecCohortId}] The number of the "extremely specific (xSpec)" cohort
definition id in the cohort table (for noisy positives).

\item[\code{daysFromxSpec}] Number of days allowed from xSpec condition until analyzed
visit

\item[\code{xSensCohortId}] The number of the "extremely sensitive (xSens)" cohort
definition id in the cohort table (for noisy negatives).

\item[\code{prevalenceCohortId}] The number of the cohort definition id to determine the
disease prevalence.

\item[\code{caseCohortId}] The number of the cohort definition id to determine cases
in the evaluation cohort

\item[\code{caseFirstOccurrenceOnly}] Set to true if only the first occurrence per subject in
the case cohort is to be used

\item[\code{xSpecCohortSize}] The recommended xSpec sample size to use in model (default
= NULL)

\item[\code{covariateSettings}] A covariateSettings object as generated using
createCovariateSettings().

\item[\code{modelPopulationCohortId}] The number of the cohort to be used as a base population
for the model. If set to 0, the entire database population
will be used.

\item[\code{modelPopulationCohortIdStartDay}] The number of days relative to the mainPopulationCohortId
cohort start date to begin including visits.

\item[\code{modelPopulationCohortIdEndDay}] The number of days relative to the mainPopulationCohortId
cohort start date to end including visits.

\item[\code{inclusionEvaluationCohortId}] The number of the cohort of the population to be used to
designate which visits are eligible to be in the
evaluation cohort

\item[\code{inclusionEvaluationDaysFromStart}] The number of days from the cohort start date of the
inclusionEvaluationCohortId to start eligible included
visits

\item[\code{inclusionEvaluationDaysFromEnd}] The number of days from the cohort start date of the
inclusionEvaluationCohortId to end eligible included
visits

\item[\code{duringInclusionEvaluationOnly}] Only include visits that are within the cohort start and
end dates

\item[\code{exclusionEvaluationCohortId}] The number of the cohort of the population to be used to
designate which visits are NOT eligible to be in the
evaluation cohort

\item[\code{exclusionEvaluationDaysFromStart}] The number of days from the cohort start date of the
exclusionEvaluationCohortId to start ineligible included
visits

\item[\code{exclusionEvaluationDaysFromEnd}] The number of days from the cohort start date of the
exclusionEvaluationCohortId to end ineligible included
visits

\item[\code{priorModelToUse}] folder where a previously developed model to use in
analysis will be found

\item[\code{minimumOffsetFromStart}] Minimum number of days to offset for the analysis visit
from the start of the observation period

\item[\code{minimumOffsetFromEnd}] Minimum number of days to offset for the analysis visit
from the end of the observation period

\item[\code{modelBaseSampleSize}] The number of non-xSpec subjects to include in the model

\item[\code{baseSampleSize}] The maximum number of subjects in the evaluation cohort.

\item[\code{lowerAgeLimit}] The lower age for subjects in the model.

\item[\code{upperAgeLimit}] The upper age for subjects in the model.

\item[\code{visitLength}] The minimum length of index visit for acute outcomes.

\item[\code{visitType}] The concept\_id for the visit type.

\item[\code{gender}] The gender(s) to be included.

\item[\code{race}] The race(s) to be included.

\item[\code{ethnicity}] The ethnicity(s) to be included.

\item[\code{startDate}] The starting date for including subjects in the model.

\item[\code{endDate}] The ending date for including subjects in the model.

\item[\code{falsePositiveNegativeSubjects}] Number of subjects to include for evaluating false
positives and negatives

\item[\code{exportFolder}] The folder where the csv output files will be written.

\item[\code{modelId}] A string used to generate the file names for this model.

\item[\code{excludeModelFromEvaluation}] Should subjects used in the model be excluded from the
evaluation cohort?

\item[\code{randomVisitTable}] Table stored in work directory with pre-selected random
visits in format of visit\_occurrence table

\item[\code{removeSubjectsWithFutureDates}] For buggy data with data in the future: ignore subjects
with dates in the future?

\item[\code{saveEvaluationCohortPlpData}] Should the large PLP file for the evaluation cohort be
saved? To be used for debugging purposes.
\end{ldescription}
\end{Arguments}
%
\begin{Details}
Create an object defining the parameter values.
\end{Details}
\HeaderA{createDefaultCovariateSettings}{Create the default covariate settings}{createDefaultCovariateSettings}
%
\begin{Description}
Create the default covariate settings for acute or chronic diseases
\end{Description}
%
\begin{Usage}
\begin{verbatim}
createDefaultCovariateSettings(
  excludedCovariateConceptIds = c(),
  includedCovariateIds = c(),
  includedCovariateConceptIds = c(),
  addDescendantsToExclude = FALSE,
  startDayWindow1 = 0,
  endDayWindow1 = 9999,
  startDayWindow2 = NULL,
  endDayWindow2 = NULL,
  startDayWindow3 = NULL,
  endDayWindow3 = NULL
)
\end{verbatim}
\end{Usage}
%
\begin{Arguments}
\begin{ldescription}
\item[\code{excludedCovariateConceptIds}] A list of conceptIds to exclude from featureExtraction.  These
should include all concept\_ids that were used to define the
xSpec model (default=NULL)

\item[\code{includedCovariateIds}] A list of covariate IDs that should be restricted to.

\item[\code{includedCovariateConceptIds}] A list of covariate concept IDs that should be restricted to.

\item[\code{addDescendantsToExclude}] Should descendants of excluded concepts also be excluded?
(default=FALSE)

\item[\code{startDayWindow1}] The day to start time window 1 for feature extraction

\item[\code{endDayWindow1}] The day to end time window 1 for feature extraction

\item[\code{startDayWindow2}] The day to start time window 2 for feature extraction

\item[\code{endDayWindow2}] The day to end time window 2 for feature extraction

\item[\code{startDayWindow3}] The day to start time window 3 for feature extraction

\item[\code{endDayWindow3}] The day to end time window 3 for feature extraction \#'
\end{ldescription}
\end{Arguments}
%
\begin{Details}
Function to create the default covariate settings for acute or chronic diseases
\end{Details}
%
\begin{Value}
A list containing one or more covariate settings objects from the FeatureExtraction package.
If only the first time window is specified, returns a list with one element.
If additional time windows are specified, returns a list with multiple elements.
\end{Value}
%
\begin{Examples}
\begin{ExampleCode}
covSettings <- createDefaultCovariateSettings(
  excludedCovariateConceptIds = c(201826, 44786627),
  addDescendantsToExclude = TRUE
)

\end{ExampleCode}
\end{Examples}
\HeaderA{createEvaluationCohort}{Create the evaluation cohort}{createEvaluationCohort}
%
\begin{Description}
Create the evaluation cohort
\end{Description}
%
\begin{Usage}
\begin{verbatim}
createEvaluationCohort(
  connectionDetails,
  oracleTempSchema = NULL,
  tempEmulationSchema = getOption("sqlRenderTempEmulationSchema"),
  phenotype,
  analysisName,
  runDateTime,
  databaseId,
  xSpecCohortId,
  daysFromxSpec = 0,
  xSensCohortId,
  prevalenceCohortId,
  caseCohortId,
  caseFirstOccurrenceOnly,
  xSpecCohortSize = 5000,
  cdmDatabaseSchema,
  cohortDatabaseSchema,
  cohortTable,
  workDatabaseSchema,
  covariateSettings = createDefaultCovariateSettings(excludedCovariateConceptIds = c(),
    addDescendantsToExclude = TRUE),
  modelPopulationCohortId = 0,
  modelPopulationCohortIdStartDay = 0,
  modelPopulationCohortIdEndDay = 0,
  inclusionEvaluationCohortId = 0,
  inclusionEvaluationDaysFromStart = 0,
  inclusionEvaluationDaysFromEnd = 0,
  duringInclusionEvaluationOnly = FALSE,
  exclusionEvaluationCohortId = 0,
  exclusionEvaluationDaysFromStart = 0,
  exclusionEvaluationDaysFromEnd = 0,
  priorModelToUse = NULL,
  minimumOffsetFromStart = 365,
  minimumOffsetFromEnd = 365,
  modelBaseSampleSize = 25000,
  baseSampleSize = 2e+06,
  lowerAgeLimit = 0,
  upperAgeLimit = 120,
  visitLength = 0,
  visitType = c(9201, 9202, 9203, 262, 581477),
  gender = c(8507, 8532),
  race = 0,
  ethnicity = 0,
  startDate = "19001010",
  endDate = "21000101",
  falsePositiveNegativeSubjects = 10,
  cdmVersion = "5",
  outFolder,
  exportFolder,
  modelId = "main",
  evaluationCohortId = "main",
  excludeModelFromEvaluation = FALSE,
  randomVisitTable = "",
  removeSubjectsWithFutureDates = TRUE,
  saveEvaluationCohortPlpData = FALSE
)
\end{verbatim}
\end{Usage}
%
\begin{Arguments}
\begin{ldescription}
\item[\code{connectionDetails}] connectionDetails created using the function
createConnectionDetails in the DatabaseConnector package.

\item[\code{oracleTempSchema}] DEPRECATED: use \code{tempEmulationSchema} instead.

\item[\code{tempEmulationSchema}] Some database platforms like Oracle and Impala do not truly support temp tables. To
emulate temp tables, provide a schema with write privileges where temp tables
can be created.

\item[\code{phenotype}] Name of the phenotype for analysis

\item[\code{analysisName}] Name of the analysis

\item[\code{runDateTime}] Starting date and time of the PheValuator run

\item[\code{databaseId}] Name of the database in the analysis

\item[\code{xSpecCohortId}] The number of the "extremely specific (xSpec)" cohort
definition id in the cohort table (for noisy positives).

\item[\code{daysFromxSpec}] Number of days allowed from xSpec condition until analyzed visit

\item[\code{xSensCohortId}] The number of the "extremely sensitive (xSens)" cohort
definition id in the cohort table (for noisy negatives).

\item[\code{prevalenceCohortId}] The number of the cohort definition id to determine the
disease prevalence.

\item[\code{caseCohortId}] The number of the cohort definition id to determine cases in the evaluation cohort

\item[\code{caseFirstOccurrenceOnly}] Set to true if only the first occurrence per subject in the case cohort is to be used

\item[\code{xSpecCohortSize}] The recommended xSpec sample size to use in model (default = NULL)

\item[\code{cdmDatabaseSchema}] The name of the database schema that contains the OMOP CDM
instance. Requires read permissions to this database. On SQL
Server, this should specifiy both the database and the
schema, so for example 'cdm\_instance.dbo'.

\item[\code{cohortDatabaseSchema}] The name of the database schema that is the location where
the cohort data used to define the at risk cohort is
available. Requires read permissions to this database.

\item[\code{cohortTable}] The tablename that contains the at risk cohort. The
expectation is cohortTable has format of COHORT table:
cohort\_concept\_id, SUBJECT\_ID, COHORT\_START\_DATE,
COHORT\_END\_DATE.

\item[\code{workDatabaseSchema}] The name of the database schema that is the location where
a table can be created and afterwards removed.
Requires write permissions to this database.

\item[\code{covariateSettings}] A covariateSettings object as generated using
createCovariateSettings().

\item[\code{modelPopulationCohortId}] The number of the cohort to be used as a base population for
the model. If set to 0, the entire database population will be
used.

\item[\code{modelPopulationCohortIdStartDay}] The number of days relative to the mainPopulationCohortId
cohort start date to begin including visits.

\item[\code{modelPopulationCohortIdEndDay}] The number of days relative to the mainPopulationCohortId
cohort start date to end including visits.

\item[\code{inclusionEvaluationCohortId}] The number of the cohort of the population to be used to designate which visits
are eligible to be in the evaluation cohort

\item[\code{inclusionEvaluationDaysFromStart}] The number of days from the cohort start date of the inclusionEvaluationCohortId
to start eligible included visits

\item[\code{inclusionEvaluationDaysFromEnd}] The number of days from the cohort start date of the inclusionEvaluationCohortId
to end eligible included visits

\item[\code{duringInclusionEvaluationOnly}] Only include visits that are within the cohort start and end dates

\item[\code{exclusionEvaluationCohortId}] The number of the cohort of the population to be used to designate which visits
are NOT eligible to be in the evaluation cohort

\item[\code{exclusionEvaluationDaysFromStart}] The number of days from the cohort start date of the exclusionEvaluationCohortId
to start ineligible included visits

\item[\code{exclusionEvaluationDaysFromEnd}] The number of days from the cohort start date of the exclusionEvaluationCohortId
to end ineligible included visits

\item[\code{priorModelToUse}] folder where a previously developed model to use in analysis will be found

\item[\code{minimumOffsetFromStart}] Minimum number of days to offset for the analysis visit from the start of the observation period

\item[\code{minimumOffsetFromEnd}] Minimum number of days to offset for the analysis visit from the end of the observation period

\item[\code{modelBaseSampleSize}] The number of non-xSpec subjects to include in the model

\item[\code{baseSampleSize}] The maximum number of subjects in the evaluation cohort.

\item[\code{lowerAgeLimit}] The lower age for subjects in the model.

\item[\code{upperAgeLimit}] The upper age for subjects in the model.

\item[\code{visitLength}] The minimum length of index visit for acute outcomes.

\item[\code{visitType}] The concept\_id for the visit type.

\item[\code{gender}] The gender(s) to be included.

\item[\code{race}] The race(s) to be included.

\item[\code{ethnicity}] The ethnicity(s) to be included.

\item[\code{startDate}] The starting date for including subjects in the model.

\item[\code{endDate}] The ending date for including subjects in the model.

\item[\code{falsePositiveNegativeSubjects}] Number of subjects to include for evaluating false positives and negatives

\item[\code{cdmVersion}] The CDM version of the database.

\item[\code{outFolder}] The folder where the output files will be written.

\item[\code{exportFolder}] The folder where the csv output files will be written.

\item[\code{modelId}] A string used to generate the file names for this model.

\item[\code{evaluationCohortId}] A string used to generate the file names for this evaluation cohort.

\item[\code{excludeModelFromEvaluation}] Should subjects used in the model be excluded from the evaluation cohort?

\item[\code{randomVisitTable}] Table stored in work directory with pre-selected random visits in format of visit\_occurrence table

\item[\code{removeSubjectsWithFutureDates}] For buggy data with data in the future: ignore subjects with
dates in the future?

\item[\code{saveEvaluationCohortPlpData}] Should the large PLP file for the evaluation cohort be saved? To be
used for debugging purposes.
\end{ldescription}
\end{Arguments}
%
\begin{Details}
Fits a diagnostic prediction model, and uses it to create an evaluation cohort with
probabilities for the health outcome of interest.
\end{Details}
%
\begin{Value}
No return value, called for side effects. Creates database tables and output files containing
the predictive model and evaluation cohort data with predicted probabilities for the health outcome.
\end{Value}
%
\begin{Examples}
\begin{ExampleCode}
## Not run: 
connectionDetails <- DatabaseConnector::createConnectionDetails(
  dbms = "postgresql",
  server = "myserver",
  user = "user",
  password = "password"
)
createEvaluationCohort(
  connectionDetails = connectionDetails,
  phenotype = "Diabetes",
  analysisName = "Main",
  runDateTime = format(Sys.time(), "%b %d %Y %X"),
  databaseId = "myDB",
  xSpecCohortId = 1,
  xSensCohortId = 2,
  prevalenceCohortId = 3,
  caseCohortId = 3,
  caseFirstOccurrenceOnly = TRUE,
  cdmDatabaseSchema = "cdm.dbo",
  cohortDatabaseSchema = "results.dbo",
  cohortTable = "cohort",
  workDatabaseSchema = "scratch.dbo",
  outFolder = tempdir(),
  exportFolder = file.path(tempdir(), "export")
)

## End(Not run)

\end{ExampleCode}
\end{Examples}
\HeaderA{createPheValuatorAnalysis}{Create a PheValuator analysis specification}{createPheValuatorAnalysis}
%
\begin{Description}
Create a PheValuator analysis specification
\end{Description}
%
\begin{Usage}
\begin{verbatim}
createPheValuatorAnalysis(
  analysisId,
  description,
  createEvaluationCohortArgs,
  testPhenotypeAlgorithmArgs
)
\end{verbatim}
\end{Usage}
%
\begin{Arguments}
\begin{ldescription}
\item[\code{analysisId}] An integer that will be used later to refer to this specific
set of analysis choices.

\item[\code{description}] A short description of the analysis.

\item[\code{createEvaluationCohortArgs}] An object representing the arguments to be used when calling
the \code{\LinkA{createCreateEvaluationCohortArgs}{createCreateEvaluationCohortArgs}} function.

\item[\code{testPhenotypeAlgorithmArgs}] Should the \code{\LinkA{createTestPhenotypeAlgorithmArgs}{createTestPhenotypeAlgorithmArgs}} function be used in this
analysis?
\end{ldescription}
\end{Arguments}
%
\begin{Details}
Create a set of analysis choices, to be used with the \code{\LinkA{runPheValuatorAnalyses}{runPheValuatorAnalyses}} function.
\end{Details}
%
\begin{Value}
A list object of class \code{"pheValuatorAnalysis"} containing the analysis specification.
Includes the analysis ID, description, and argument specifications for creating the evaluation cohort
and testing the phenotype algorithm.
\end{Value}
%
\begin{Examples}
\begin{ExampleCode}
evalCohortArgs <- createCreateEvaluationCohortArgs(
  xSpecCohortId = 1,
  xSensCohortId = 2,
  prevalenceCohortId = 3
)
testArgs <- createTestPhenotypeAlgorithmArgs(
  phenotypeCohortId = 10
)
analysis <- createPheValuatorAnalysis(
  analysisId = 1,
  description = "Main analysis",
  createEvaluationCohortArgs = evalCohortArgs,
  testPhenotypeAlgorithmArgs = testArgs
)

\end{ExampleCode}
\end{Examples}
\HeaderA{createTestPhenotypeAlgorithmArgs}{Create a parameter object for the function testPhenotypeAlgorithm}{createTestPhenotypeAlgorithmArgs}
%
\begin{Description}
Create a parameter object for the function testPhenotypeAlgorithm
\end{Description}
%
\begin{Usage}
\begin{verbatim}
createTestPhenotypeAlgorithmArgs(
  phenotype,
  analysisName,
  runDateTime,
  cutPoints = c("EV"),
  exportFolder,
  phenotypeCohortId,
  databaseId,
  washoutPeriod = 0,
  splayPrior = 7,
  splayPost = 7
)
\end{verbatim}
\end{Usage}
%
\begin{Arguments}
\begin{ldescription}
\item[\code{phenotype}] Name of the phenotype for analysis

\item[\code{analysisName}] Name of the analysis

\item[\code{runDateTime}] Starting date and time of the PheValuator run

\item[\code{cutPoints}] A list of threshold predictions for the evaluations.  Include "EV" for
the expected value

\item[\code{exportFolder}] The folder where the csv output files will be written.

\item[\code{phenotypeCohortId}] The ID of the cohort to evaluate in the specified cohort table.

\item[\code{databaseId}] Name of the database in the analysis

\item[\code{washoutPeriod}] The mininum required continuous observation time prior to index date for
subjects within the cohort to test (Default = 0).

\item[\code{splayPrior}] The number of days to allow for test phenotype visit date prior to
evaluation date

\item[\code{splayPost}] The number of days to allow for test phenotype visit date after
evaluation date
\end{ldescription}
\end{Arguments}
%
\begin{Details}
Create an object defining the parameter values.
\end{Details}
\HeaderA{loadPheValuatorAnalysisList}{Load a list of phevaluatorAnalysis from file}{loadPheValuatorAnalysisList}
%
\begin{Description}
Load a list of objects of type \code{pheValuatorAnalysis} from file. The file is in JSON format.
\end{Description}
%
\begin{Usage}
\begin{verbatim}
loadPheValuatorAnalysisList(file)
\end{verbatim}
\end{Usage}
%
\begin{Arguments}
\begin{ldescription}
\item[\code{file}] The name of the file
\end{ldescription}
\end{Arguments}
%
\begin{Value}
A list of objects of type \code{pheValuatorAnalysis}.
\end{Value}
%
\begin{Examples}
\begin{ExampleCode}
## Not run: 
analysisList <- loadPheValuatorAnalysisList(file = "analysisList.json")

## End(Not run)

\end{ExampleCode}
\end{Examples}
\HeaderA{runPheValuatorAnalyses}{Run a list of analyses}{runPheValuatorAnalyses}
%
\begin{Description}
Run a list of analyses
\end{Description}
%
\begin{Usage}
\begin{verbatim}
runPheValuatorAnalyses(
  phenotype,
  cohortDefinitionSet = data.frame(),
  analysisName = "Main",
  runDateTime = format(Sys.time(), "%b %d %Y %X"),
  connectionDetails,
  oracleTempSchema = NULL,
  tempEmulationSchema = getOption("sqlRenderTempEmulationSchema"),
  cdmDatabaseSchema,
  cohortDatabaseSchema = cdmDatabaseSchema,
  cohortTable = "cohort",
  workDatabaseSchema = cdmDatabaseSchema,
  databaseId = cdmDatabaseSchema,
  cdmVersion = 5,
  outputFolder,
  priorModelToUse = NULL,
  pheValuatorAnalysisList
)
\end{verbatim}
\end{Usage}
%
\begin{Arguments}
\begin{ldescription}
\item[\code{phenotype}] Name of the phenotype for analysis

\item[\code{cohortDefinitionSet}] Data.frame of cohorts must include columns cohortId, cohortName, json, sql - should include all
cohort definitions needed to replicate PheValuator analysis

\item[\code{analysisName}] Name of the analysis

\item[\code{runDateTime}] Starting date and time of the PheValuator run

\item[\code{connectionDetails}] An R object of type \code{connectionDetails} created using the
function \code{createConnectionDetails} in the
\code{DatabaseConnector} package.

\item[\code{oracleTempSchema}] DEPRECATED: use \code{tempEmulationSchema} instead.

\item[\code{tempEmulationSchema}] Some database platforms like Oracle and Impala do not truly support temp tables. To
emulate temp tables, provide a schema with write privileges where temp tables
can be created.

\item[\code{cdmDatabaseSchema}] The name of the database schema that contains the OMOP CDM
instance. Requires read permissions to this database. On SQL
Server, this should specifiy both the database and the
schema, so for example 'cdm\_instance.dbo'.

\item[\code{cohortDatabaseSchema}] The name of the database schema that is the location where
the cohort data used to define the at risk cohort is
available. Requires read permissions to this database.

\item[\code{cohortTable}] The tablename that contains the at risk cohort. The
expectation is cohortTable has format of COHORT table:
cohort\_concept\_id, SUBJECT\_ID, COHORT\_START\_DATE,
COHORT\_END\_DATE.

\item[\code{workDatabaseSchema}] The name of the database schema that is the location where
a table can be created and afterwards removed.
Requires write permissions to this database.

\item[\code{databaseId}] Name of the database in the analysis

\item[\code{cdmVersion}] Define the OMOP CDM version used: currently supports "5".

\item[\code{outputFolder}] Name of the folder where all the outputs will be written to.

\item[\code{priorModelToUse}] folder where a previously developed model to use in analysis will be found

\item[\code{pheValuatorAnalysisList}] A list of objects of type \code{pheValuatorAnalysis} as created using
the \code{\LinkA{createPheValuatorAnalysis}{createPheValuatorAnalysis}} function.
\end{ldescription}
\end{Arguments}
%
\begin{Details}
Run a list of analyses.
\end{Details}
%
\begin{Value}
A data frame specifying where the constructed evaluation cohort and phenotype evaluation results can be found
in the local file system.
\end{Value}
%
\begin{Examples}
\begin{ExampleCode}
## Not run: 
connectionDetails <- DatabaseConnector::createConnectionDetails(
  dbms = "postgresql",
  server = "myserver",
  user = "user",
  password = "password"
)
analysisList <- list(
  createPheValuatorAnalysis(
    analysisId = 1,
    description = "Main analysis",
    createEvaluationCohortArgs = createCreateEvaluationCohortArgs(
      xSpecCohortId = 1,
      prevalenceCohortId = 3
    ),
    testPhenotypeAlgorithmArgs = createTestPhenotypeAlgorithmArgs(
      phenotypeCohortId = 10
    )
  )
)
runPheValuatorAnalyses(
  phenotype = "Diabetes",
  connectionDetails = connectionDetails,
  cdmDatabaseSchema = "cdm.dbo",
  outputFolder = tempdir(),
  pheValuatorAnalysisList = analysisList
)

## End(Not run)

\end{ExampleCode}
\end{Examples}
\HeaderA{savePheValuatorAnalysisList}{Save a list of phevaluatorAnalysis to file}{savePheValuatorAnalysisList}
%
\begin{Description}
Write a list of objects of type \code{pheValuatorAnalysis} to file. The file is in JSON format.
\end{Description}
%
\begin{Usage}
\begin{verbatim}
savePheValuatorAnalysisList(pheValuatorAnalysisList, file)
\end{verbatim}
\end{Usage}
%
\begin{Arguments}
\begin{ldescription}
\item[\code{pheValuatorAnalysisList}] The pheValuatorAnalysis list to be written to file

\item[\code{file}] The name of the file where the results will be written
\end{ldescription}
\end{Arguments}
%
\begin{Value}
No return value, called for side effects. Writes a JSON file containing the analysis list
to the specified file path.
\end{Value}
%
\begin{Examples}
\begin{ExampleCode}
## Not run: 
# Requires pre-existing analysis list object; example usage:
# savePheValuatorAnalysisList(pheValuatorAnalysisList = myAnalysisList,
#                             file = file.path(tempdir(), "analysisList.json"))

## End(Not run)

\end{ExampleCode}
\end{Examples}
\HeaderA{summarizePheValuatorAnalyses}{Summarize results of PheValuator analyses}{summarizePheValuatorAnalyses}
%
\begin{Description}
Summarize results of PheValuator analyses
\end{Description}
%
\begin{Usage}
\begin{verbatim}
summarizePheValuatorAnalyses(referenceTable, outputFolder)
\end{verbatim}
\end{Usage}
%
\begin{Arguments}
\begin{ldescription}
\item[\code{referenceTable}] A reference table as created using \code{\LinkA{runPheValuatorAnalyses}{runPheValuatorAnalyses}}.

\item[\code{outputFolder}] The output folder used when calling \code{\LinkA{runPheValuatorAnalyses}{runPheValuatorAnalyses}}.
\end{ldescription}
\end{Arguments}
%
\begin{Value}
A data frame of results.
\end{Value}
%
\begin{Examples}
\begin{ExampleCode}
## Not run: 
refTable <- runPheValuatorAnalyses(
  phenotype = "Diabetes",
  connectionDetails = connectionDetails,
  cdmDatabaseSchema = "cdm.dbo",
  outputFolder = tempdir(),
  pheValuatorAnalysisList = analysisList
)
results <- summarizePheValuatorAnalyses(refTable, outputFolder = tempdir())

## End(Not run)

\end{ExampleCode}
\end{Examples}
\HeaderA{testPhenotypeAlgorithm}{Test phenotype algorithms}{testPhenotypeAlgorithm}
%
\begin{Description}
Test phenotype algorithms
\end{Description}
%
\begin{Usage}
\begin{verbatim}
testPhenotypeAlgorithm(
  phenotype,
  analysisName,
  runDateTime,
  connectionDetails,
  cutPoints = c("EV"),
  outFolder,
  exportFolder,
  evaluationCohortId = "main",
  phenotypeCohortId,
  cdmDatabaseSchema,
  databaseId,
  cohortDatabaseSchema,
  cohortTable,
  washoutPeriod = 0,
  splayPrior = 7,
  splayPost = 7
)
\end{verbatim}
\end{Usage}
%
\begin{Arguments}
\begin{ldescription}
\item[\code{phenotype}] Name of the phenotype for analysis

\item[\code{analysisName}] Name of the analysis

\item[\code{runDateTime}] Starting date and time of the PheValuator run

\item[\code{connectionDetails}] ConnectionDetails created using the function
createConnectionDetails in the DatabaseConnector package.

\item[\code{cutPoints}] A list of threshold predictions for the evaluations.  Include "EV"
for the expected value

\item[\code{outFolder}] The folder where the cohort evaluation output files are written

\item[\code{exportFolder}] The folder where the csv output files will be written.

\item[\code{evaluationCohortId}] A string used to generate the file names for the evaluation cohort.

\item[\code{phenotypeCohortId}] The ID of the cohort to evaluate in the specified cohort table.

\item[\code{cdmDatabaseSchema}] The name of the database schema that contains the OMOP CDM
instance. Requires read permissions to this database. On SQL
Server, this should specifiy both the database and the
schema, so for example 'cdm\_instance.dbo'.

\item[\code{databaseId}] Name of the database in the analysis

\item[\code{cohortDatabaseSchema}] The name of the database schema that is the location where the
cohort data used to define the at risk cohort is available.
Requires read permissions to this database.

\item[\code{cohortTable}] The tablename that contains the at risk cohort. The expectation is
cohortTable has format of COHORT table: cohort\_concept\_id,
SUBJECT\_ID, COHORT\_START\_DATE, COHORT\_END\_DATE.

\item[\code{washoutPeriod}] The mininum required continuous observation time prior to index
date for subjects within the cohort to test (Default = 0).

\item[\code{splayPrior}] The number of days to allow for test phenotype visit date prior to evaluation date

\item[\code{splayPost}] The number of days to allow for test phenotype visit date after evaluation date
\end{ldescription}
\end{Arguments}
%
\begin{Details}
This function will perform the phenotype algorithm evaluation using the evaluation cohort returned
from createEvalCohort and the phenotype algorithm cohort specified
\end{Details}
%
\begin{Value}
A dataframe with the results from the phenotype
algorithm evaluation.

If 0.5 is included as a cutpoint, the data frame will have an attribute called 'misses', a dataframe
with a sample of subject ids for TPs, FPs, TNs, and FNs for the 50 percent and over prediction threshold.
\end{Value}
%
\begin{Examples}
\begin{ExampleCode}
## Not run: 
connectionDetails <- DatabaseConnector::createConnectionDetails(
  dbms = "postgresql",
  server = "myserver",
  user = "user",
  password = "password"
)
results <- testPhenotypeAlgorithm(
  phenotype = "Diabetes",
  analysisName = "Main",
  runDateTime = format(Sys.time(), "%b %d %Y %X"),
  connectionDetails = connectionDetails,
  cutPoints = c("EV", "0.5"),
  outFolder = tempdir(),
  exportFolder = file.path(tempdir(), "export"),
  phenotypeCohortId = 10,
  cdmDatabaseSchema = "cdm.dbo",
  databaseId = "myDB",
  cohortDatabaseSchema = "results.dbo",
  cohortTable = "cohort"
)

## End(Not run)

\end{ExampleCode}
\end{Examples}
\printindex{}
\end{document}
